% ============================================================
%  The European AI Summer Research
%  Paper template (recommended, optional)
%
%  HOW TO USE
%  ----------
%  1. Open this folder on Overleaf (overleaf.com) or compile
%     locally with pdflatex. Keep easrp2026.sty in the same folder.
%  2. Write your paper in the marked sections below.
%  3. Keep it ANONYMOUS: no names, no affiliations, no
%     identifying acknowledgements. See README.txt.
%  4. Compile to PDF and submit the PDF.
%
%  PAGE LIMIT: up to 8 pages of main text (11pt, A4).
%  References and any appendix do NOT count toward this limit.
%
%  When your paper is ACCEPTED, change the line
%      \usepackage{easrp2026}
%  to
%      \usepackage[accepted]{easrp2026}
%  and fill in \easrpauthor / \easrpaffiliation below to show
%  your name on the final version.
% ============================================================

\documentclass[11pt,a4paper]{article}

\usepackage{graphicx}             % for inserting images
\usepackage{easrp2026}            % anonymous submission
% \usepackage[accepted]{easrp2026}  % camera-ready (reveals authors)

% --- Author details: used only in the accepted version -------
% Leave these here; they are ignored while in anonymous mode.
\easrpauthor{Your Name, Co-author Name}
\easrpaffiliation{Your University, City, Country}

\easrptitle{Your paper title here}

\begin{document}

\maketitle

% Delete this note in your final paper.
\begin{center}
\fbox{\parbox{0.9\linewidth}{\small\itshape
Formatting: up to 8 pages of main text at 11pt on A4; references
and any appendix do not count toward the limit. Keep the
submission anonymous (no names, affiliations, or identifying
acknowledgements).}}
\end{center}

\begin{abstract}
Write a short abstract here, roughly 150--250 words: the
problem, your approach, and your main result, in plain
language. Do not include anything that identifies you or your
institution.
\end{abstract}

\section{Introduction}
Introduce the problem and why it matters, and state your
contribution clearly. Write in an anonymous style: refer to
your own prior work in the third person (for example,
``Smith (2024) showed\ldots'') rather than ``our previous
work''.

\section{Related work}
Situate your work among existing research: summarise what has
been done, identify what remains open, and make clear how your
contribution differs. Cite credible, verifiable sources, and
prefer peer-reviewed papers where they exist. Preprints,
technical reports, datasets, theses, and software may be cited
when they are the appropriate source (as is common in machine
learning); avoid relying on non-authoritative web content such
as blog posts.

The accompanying bibliography file includes examples of several
reference types: a journal article \cite{lecun2015deep}, a
conference paper \cite{vaswani2017attention}, a book
\cite{russell2020aima}, a book chapter \cite{rumelhart1986learning},
a technical report \cite{krizhevsky2009cifar}, a thesis
\cite{sutton1984thesis}, a preprint \cite{devlin2018bert}, and a
web resource \cite{turing1950webexample}.

This template uses \texttt{natbib}, which offers several citation
commands. Use \verb|\citet| for a textual citation that reads as
part of the sentence: \citet{vaswani2017attention} introduced the
Transformer. Use \verb|\citep| for a parenthetical citation: the
Transformer dispenses with recurrence \citep{vaswani2017attention}.
The starred forms \verb|\citet*| and \verb|\citep*| list every
author rather than abbreviating with ``et al.'':
\citep*{vaswani2017attention}. You can also cite just the year with
\verb|\citeyear| \citeyear{lecun2015deep}, just the authors with
\verb|\citeauthor| \citeauthor{lecun2015deep}, or give multiple
references at once \citep{lecun2015deep, devlin2018bert}.

\section{Method}
Describe your approach, model, or argument with enough detail
that an expert could follow and, where applicable, reproduce
it.

\section{Results}
Present your findings. A rigorous negative result is valued:
you will not be penalised for failing to beat the state of the
art.

\section{Discussion}
Interpret your results, state limitations honestly, and
suggest future directions.

\section{Conclusion}
Summarise your contribution in a few sentences.

% --- References ----------------------------------------------
% References and any appendix do NOT count toward the 8-page limit.
\bibliographystyle{plainnat}
\bibliography{references}

% --- Appendix (optional; does NOT count toward the limit) ----
\appendix
\section{Appendix}
Supplementary material, proofs, or extra detail may go here.

\end{document}